%%%%%%%%%%%%%%%%%%%%%%%%%%%%%%%%%%%%%%%%%%%%%%%%%%%%%%%%%%%%%%%%%
% PHIẾU HỌC TẬP: ECOFOOTPRINT
%%%%%%%%%%%%%%%%%%%%%%%%%%%%%%%%%%%%%%%%%%%%%%%%%%%%%%%%%%%%%%%%%
\documentclass[12pt,a4paper]{article}

% ============= GÓI CƠ BẢN =============
\usepackage[utf8]{inputenc}
\usepackage[T1]{fontenc}
\usepackage[vietnamese]{babel}
\usepackage{graphicx}
\usepackage{amsmath,amssymb}
\usepackage{tabularx}
\usepackage{colortbl}
\usepackage{tikz}
\usepackage{enumitem}
\usepackage{geometry}
\geometry{left=2cm,right=2cm,top=2cm,bottom=2cm}

% ============= MÀU SẮC =============
\usepackage[svgnames]{xcolor}
\definecolor{primary}{RGB}{34,139,34}
\definecolor{headerbg}{RGB}{200,240,200}

% ============= TIÊU ĐỀ =============
\newcommand{\MyHeader}{
    \begin{tikzpicture}[remember picture, overlay]
        \draw[primary, line width=2pt] ($(current page.north west)+(1.5,-1.5)$) rectangle ($(current page.south east)+(-1.5,1.5)$);
    \end{tikzpicture}
    \begin{center}
        {\large\sffamily SỞ GD\&ĐT TP. HỒ CHÍ MINH}\\
        {\bfseries\large\sffamily TRƯỜNG THPT NGUYỄN HỮU CẢNH}\\[0.5cm]
        {\huge\bfseries\color{primary}\sffamily PHIẾU HỌC TẬP STEM}\\[0.2cm]
        {\Large\sffamily CHỦ ĐỀ: TÍNH DẤU CHÂN CARBON CÁ NHÂN}\\[0.5cm]
    \end{center}
    \noindent\textbf{Họ và tên:} \underline{\hspace{6cm}} \textbf{Lớp:} \underline{\hspace{2cm}} \textbf{Nhóm:} \underline{\hspace{1cm}}
    \vspace{0.5cm}
    \hrule
    \vspace{0.5cm}
}

\begin{document}
\MyHeader

% PHẦN 1
\section*{\color{primary} I. KHẢO SÁT DỮ LIỆU ĐẦU VÀO (DATA COLLECTION)}
\textit{Hãy thu thập số liệu thực tế từ gia đình em trong 1 tháng vừa qua.}

\begin{center}
\renewcommand{\arraystretch}{1.5}
\begin{tabularx}{\textwidth}{|>{\columncolor{headerbg}\bfseries}l|X|X|}
\hline
\textbf{Nguồn phát thải} & \textbf{Số liệu thu thập} & \textbf{Quy đổi về đơn vị chuẩn} \\
\hline
\textbf{1. Điện năng} & Hóa đơn tháng trước: \dotfill VND & Số kWh tiêu thụ = $\dfrac{\text{Tiền}}{2.500} \approx$ \dotfill kWh \\
\hline
\textbf{2. Nước sạch} & Hóa đơn nước: \dotfill VND & Số $m^3$ tiêu thụ = $\dfrac{\text{Tiền}}{10.000} \approx$ \dotfill $m^3$ \\
\hline
\textbf{3. Di chuyển} & Phương tiện chính: \dotfill & Quãng đường đi/ngày: \dotfill km \\
(Xe máy/Bus/...) & Số ngày đi/tuần: \dotfill & Tổng km/tháng $\approx$ \dotfill km \\
\hline
\textbf{4. Rác thải} & Ước lượng rác nhà em: \dotfill kg/ngày & Tổng kg/tháng $\approx$ \dotfill kg \\
\hline
\end{tabularx}
\end{center}

% PHẦN 2
\section*{\color{primary} II. TÍNH TOÁN TOÁN HỌC (MATHEMATICAL MODELING)}
\textbf{Công thức:} $E = EF \times x$ (trong đó $EF$ là hệ số phát thải).

\begin{enumerate}
    \item \textbf{Tính lượng $CO_2$ từ ĐIỆN:} ($EF_{dien} = 0.72$ kg/kWh)
    $$E_{dien} = 0.72 \times \text{\hspace{2cm}} = \text{\hspace{3cm}} \text{ (kg } CO_2)$$
    
    \item \textbf{Tính lượng $CO_2$ từ XE MÁY:} ($EF_{xe} = 0.05$ kg/km)
    $$E_{xe} = 0.05 \times \text{\hspace{2cm}} = \text{\hspace{3cm}} \text{ (kg } CO_2)$$
    
    \item \textbf{Tổng Dấu chân Carbon ước tính trong 1 tháng:}
    $$E_{total} = E_{dien} + E_{xe} + E_{nuoc} + \dots = \text{\hspace{3cm}} \text{ (kg } CO_2)$$
    
    \item \textbf{Quy đổi ra 1 năm:} $E_{nam} = E_{total} \times 12 = \text{\hspace{3cm}} \text{ (kg } CO_2)$
\end{enumerate}

% PHẦN 3
\section*{\color{primary} III. ỨNG DỤNG CÔNG NGHỆ (ECO APP)}
Truy cập: \textbf{\texttt{stem-1h8.pages.dev/eco.html}}

\begin{enumerate}
    \item Nhập các số liệu em đã thu thập vào ứng dụng.
    \item Ghi lại kết quả từ App:
    \begin{itemize}
        \item Tổng phát thải: \dotfill kg/năm. (So sánh với tính tay: lệch bao nhiêu? \dotfill)
        \item Nguồn phát thải lớn nhất là gì? \dotfill (chiếm \dotfill \%)
    \end{itemize}
    \item \textbf{Thử nghiệm kịch bản giảm thải:} 
    \textit{Nếu em giảm dùng điện 10\% và đi xe đạp 2 ngày/tuần, App dự báo em sẽ giảm được bao nhiêu kg $CO_2$?}
    \vspace{1cm}
\end{enumerate}

% PHẦN 4
\section*{\color{primary} IV. CAM KẾT HÀNH ĐỘNG}
\textit{Dựa trên kết quả tính toán, hãy viết ra 3 hành động cụ thể em sẽ làm:}
\begin{itemize}
    \item [1.] \dotfill
    \item [2.] \dotfill
    \item [3.] \dotfill
\end{itemize}

\end{document}
